\section{Arquitectura de blockchain}
\label{sec:arquitectura_blockchain}

Una \textit{blockchain} es un registro distribuido (\textit{distributed ledger}) de
transacciones que se mantiene de forma replicada entre los nodos de una red \textit{peer-to-peer},
sin depender de una autoridad central que valide o custodie su contenido. Su propiedad
fundamental es que, una vez que un dato se incorpora al registro, modificarlo o eliminarlo
requiere rehacer todo el trabajo criptográfico posterior a ese punto, lo cual se vuelve
computacionalmente inviable a medida que la cadena crece. Esta sección describe los cuatro
componentes que en conjunto sostienen esa propiedad: la estructura de bloques encadenados por
hash, los árboles de Merkle que resumen su contenido, los mecanismos de consenso que deciden
qué bloque es válido, y los roles que los distintos tipos de nodos cumplen dentro de la red.
Estos conceptos son la base sobre la que se apoyan tanto el diseño de Ethereum
(Sección~\ref{sec:ethereum}) como el esquema de verificación de conocimiento cero desarrollado
en este trabajo.

\subsection{Estructura de bloques y encadenamiento por hash}

\begin{definition}
	Un \textbf{bloque} es una estructura de datos compuesta por un \textit{encabezado}
	(\textit{header}) y un cuerpo. El encabezado contiene, como mínimo, el hash criptográfico
	del encabezado del bloque anterior, una marca de tiempo, un valor que resume el conjunto de
	transacciones del bloque (la raíz de un árbol de Merkle, ver Definición~\ref{def:merkle}), y
	metadatos propios del mecanismo de consenso utilizado (por ejemplo, un \textit{nonce} en
	\textit{Proof-of-Work}). El cuerpo contiene la lista de transacciones incluidas en el bloque.
\end{definition}

La \textbf{blockchain} propiamente dicha es la secuencia de bloques $B_0, B_1, \dots, B_n$
donde $B_0$ es el \textit{bloque génesis} y cada bloque $B_i$, con $i > 0$, referencia a su
predecesor incluyendo $H(B_{i-1})$ en su propio encabezado, siendo $H$ una función de hash
criptográfica (típicamente SHA-256 o Keccak-256). Esta construcción encadenada es lo que le da
su nombre a la estructura y es la responsable directa de su inmutabilidad: como el hash de un
bloque depende del hash del bloque anterior, alterar el contenido de $B_i$ cambia su hash
$H(B_i)$, lo que invalida la referencia almacenada en $B_{i+1}$, y por extensión la de todos
los bloques posteriores hasta la cabeza de la cadena. Un atacante que quisiera modificar una
transacción confirmada tendría entonces que recalcular el trabajo de consenso de todos los
bloques desde ese punto en adelante, y hacerlo más rápido que el resto de la red honesta le
suma nuevos bloques a la cadena legítima.

\subsection{Árboles de Merkle y pruebas de inclusión}
\label{sec:merkle}

Incluir el hash de cada transacción individual en el encabezado del bloque sería costoso en
espacio y haría que verificar la pertenencia de una transacción específica implicara descargar
el bloque completo. Los árboles de Merkle~\cite{merkle1988digital} resuelven ambos problemas.

\begin{definition}
	\label{def:merkle}
	Un \textbf{árbol de Merkle} es un árbol binario en el que cada hoja contiene el hash de un
	elemento de datos (en este contexto, una transacción) y cada nodo interno contiene el hash
	de la concatenación de sus dos hijos: $h_{padre} = H(h_{izq} \Vert h_{der})$. El hash
	almacenado en la raíz, llamado \textbf{raíz de Merkle}, es un compromiso criptográfico sobre
	el conjunto completo de hojas: cambiar cualquier transacción altera su hash hoja y, por
	propagación, la raíz.
\end{definition}

Esta estructura permite generar una \textbf{prueba de inclusión} (\textit{Merkle proof}) de que
una transacción $t$ forma parte de un bloque sin necesitar el bloque completo: basta con
entregar $t$ junto con los hashes hermanos de cada nivel del camino desde su hoja hasta la
raíz. Quien verifica recalcula los hashes de ese camino y comprueba que el resultado coincide
con la raíz de Merkle publicada en el encabezado del bloque. Para un árbol con $n$ hojas, la
prueba tiene tamaño $O(\log n)$, en contraste con el $O(n)$ que implicaría transmitir todas las
transacciones. Ethereum extiende esta idea con \textit{Merkle-Patricia tries}, que además de
las transacciones permiten probar el estado de una cuenta o el resultado de una ejecución
(Sección~\ref{sec:ethereum}); esta capacidad de generar pruebas de inclusión compactas y
verificables es, precisamente, el mecanismo que este trabajo utiliza para certificar
\textit{on-chain}, mediante una prueba de conocimiento cero, que una transacción de bloqueo fue
efectivamente incluida en la cadena de origen.

\subsection{Mecanismos de consenso}

Que cada nodo mantenga una copia del registro no es suficiente: en una red sin autoridad
central, distintos nodos pueden recibir transacciones en distinto orden o proponer bloques en
simultáneo, y es necesario que todos converjan a una única secuencia de bloques válida. Este es
un caso particular del problema de los generales bizantinos~\cite{lamport1982byzantine}: lograr
acuerdo entre participantes distribuidos cuando una fracción de ellos puede ser maliciosa o
fallar arbitrariamente. Un \textbf{mecanismo de consenso} es el protocolo que resuelve este
problema definiendo, ante bloques o cadenas en conflicto, cuál es la versión canónica.

Los dos mecanismos dominantes son:

\begin{itemize}
	\item \textbf{Proof-of-Work (PoW)}: introducido por Bitcoin~\cite{nakamoto2008bitcoin}, exige
	que quien propone un bloque resuelva un problema computacional costoso (encontrar un
	\textit{nonce} tal que $H(\text{header})$ sea menor a un umbral de dificultad). La cadena
	válida es la que acumula mayor trabajo computacional total, lo que hace que reescribir el
	historial requiera superar la capacidad de cómputo de toda la red honesta. Su finalidad es
	\textit{probabilística}: la confianza en que un bloque no será revertido crece con la
	cantidad de bloques confirmados sobre él, pero nunca es absoluta.

	\item \textbf{Proof-of-Stake (PoS)}: sustituye el costo computacional por un costo
	económico. Los participantes (\textit{validadores}) bloquean una cantidad de la
	criptomoneda nativa como garantía y son seleccionados para proponer y atestiguar bloques en
	proporción a su participación bloqueada (\textit{stake}). Un comportamiento malicioso
	detectable (como firmar dos bloques en conflicto para la misma posición) resulta en la
	pérdida de parte o la totalidad de esa garantía (\textit{slashing}). A diferencia de PoW,
	las variantes modernas de PoS —como la que adoptó Ethereum tras \textit{The Merge}— pueden
	alcanzar \textbf{finalidad determinística}: una vez que un bloque acumula suficientes
	atestiguaciones, revertirlo es no solo improbable sino que exige destruir económicamente a
	los validadores responsables. Este mecanismo de finalidad, y el comité que la certifica en
	Ethereum, se retoma en la Sección~\ref{sec:ethereum} y es la raíz de confianza sobre la que
	se construye el puente propuesto en este trabajo.
\end{itemize}

\subsection{Roles de nodos y topología de red}

No todos los participantes de una red blockchain cumplen la misma función ni requieren los
mismos recursos:

\begin{itemize}
	\item \textbf{Nodos completos} (\textit{full nodes}): descargan y validan de forma
	independiente cada bloque y cada transacción desde el bloque génesis, manteniendo una copia
	íntegra del estado de la red. No confían en ningún tercero: verifican por sí mismos que se
	respetaron las reglas de consenso y de validez de transacciones.

	\item \textbf{Mineros o validadores}: además de validar, participan activamente en la
	producción de nuevos bloques (resolviendo el rompecabezas de PoW o siendo seleccionados y
	atestiguando en PoS).

	\item \textbf{Clientes ligeros} (\textit{light clients}): no almacenan la cadena completa ni
	reejecutan todas las transacciones; en su lugar, confían en los encabezados de bloque —en
	particular en su raíz de Merkle y, en redes PoS, en la firma del comité que certifica su
	finalidad— para verificar de forma eficiente si una transacción específica fue incluida,
	mediante una prueba de inclusión como la descrita en la Sección~\ref{sec:merkle}. Este es
	precisamente el modelo de confianza que adopta el \textit{relayer} \textit{off-chain} del
	puente diseñado en este trabajo: no re-ejecuta la cadena de origen, sino que verifica una
	prueba de inclusión respaldada por la firma del comité de sincronización de Ethereum.
\end{itemize}

Todos estos roles se comunican mediante una red \textit{peer-to-peer} de propagación por
\textit{gossip}, en la que cada nodo retransmite a sus pares los bloques y transacciones que
recibe y considera válidos, sin que exista un servidor central que coordine la difusión de la
información.
